% Options for packages loaded elsewhere
\PassOptionsToPackage{unicode}{hyperref}
\PassOptionsToPackage{hyphens}{url}
\documentclass[
  ukrainian,
]{article}
\usepackage{xcolor}
\usepackage[margin=2.5cm]{geometry}
\usepackage{amsmath,amssymb}
\setcounter{secnumdepth}{-\maxdimen} % remove section numbering
\usepackage{iftex}
\ifPDFTeX
  \usepackage[T1]{fontenc}
  \usepackage[utf8]{inputenc}
  \usepackage{textcomp} % provide euro and other symbols
\else % if luatex or xetex
  \usepackage{unicode-math} % this also loads fontspec
  \defaultfontfeatures{Scale=MatchLowercase}
  \defaultfontfeatures[\rmfamily]{Ligatures=TeX,Scale=1}
\fi
\usepackage{lmodern}
\ifPDFTeX\else
\setmainfont{DejaVu Serif}
\setsansfont{DejaVu Sans}
\setmonofont{DejaVu Sans Mono}
\setmathfont{STIX Two Math}
\fi
\ifPDFTeX\else
  % xetex/luatex font selection
\fi
% Use upquote if available, for straight quotes in verbatim environments
\IfFileExists{upquote.sty}{\usepackage{upquote}}{}
\IfFileExists{microtype.sty}{% use microtype if available
  \usepackage[]{microtype}
  \UseMicrotypeSet[protrusion]{basicmath} % disable protrusion for tt fonts
}{}
\makeatletter
\@ifundefined{KOMAClassName}{% if non-KOMA class
  \IfFileExists{parskip.sty}{%
    \usepackage{parskip}
  }{% else
    \setlength{\parindent}{0pt}
    \setlength{\parskip}{6pt plus 2pt minus 1pt}}
}{% if KOMA class
  \KOMAoptions{parskip=half}}
\makeatother
\usepackage{graphicx}
\graphicspath{{/Users/kalinicenkovitalijmikolajovic/.codex/worktrees/5c15/Script_for_translate_book/Output_Final/}}
\setkeys{Gin}{width=0.82\linewidth,height=0.82\textheight,keepaspectratio}
\makeatletter
\newsavebox\pandoc@box
\newcommand*\pandocbounded[1]{% scales image to fit in text height/width
  \sbox\pandoc@box{#1}%
  \Gscale@div\@tempa{\textheight}{\dimexpr\ht\pandoc@box+\dp\pandoc@box\relax}%
  \Gscale@div\@tempb{\linewidth}{\wd\pandoc@box}%
  \ifdim\@tempb\p@<\@tempa\p@\let\@tempa\@tempb\fi% select the smaller of both
  \ifdim\@tempa\p@<\p@\scalebox{\@tempa}{\usebox\pandoc@box}%
  \else\usebox{\pandoc@box}%
  \fi%
}
% Set default figure placement to htbp
\def\fps@figure{htbp}
\makeatother
\ifLuaTeX
\usepackage[bidi=basic,shorthands=off,provide=*]{babel}
\else
\usepackage[bidi=default,shorthands=off,provide=*]{babel}
\fi
\ifLuaTeX
  \usepackage{selnolig} % disable illegal ligatures
\fi
\setlength{\emergencystretch}{3em} % prevent overfull lines
\providecommand{\tightlist}{%
  \setlength{\itemsep}{0pt}\setlength{\parskip}{0pt}}
\usepackage{bookmark}
\IfFileExists{xurl.sty}{\usepackage{xurl}}{} % add URL line breaks if available
\urlstyle{same}
\hypersetup{
  pdflang={uk},
  hidelinks,
  pdfcreator={LaTeX via pandoc}}

\author{}
\date{}

\begin{document}

\section{Модель колеблющейся
балки}\label{ux43cux43eux434ux435ux43bux44c-ux43aux43eux43bux435ux431ux43bux44eux449ux435ux439ux441ux44f-ux431ux430ux43bux43aux438}

Рассмотрим свободно опирающуюся балку, изображенную на рис. 21.2. Пусть
начальные условия таковы:

\[
C_n = \frac{2}{\pi} \int_0^{\pi} f(x) \cos(nx) dx.
\]

\(u_1(x, 0) = 0\) Мы могли бы найти решение, подставив функции \(f(x)\)
и \(g(x)\) в формулы (21.6), однако легко видеть, что

\[
y = \frac{\pi}{b-a}(x-a).
\]

для всех \(n=1,\,2,\,\ldots,\) \(b_1=1,\,b_2=0,\,b_3=0,5\)
\(b_n=0,\,n=4,\,5,\,\ldots\)

Следовательно, решение имеет вид

\[
f(x) = \frac{C_0}{2} + \sum_{n=1}^{\infty} C_n \cos(nx).
\]

.

Интересно сравнить это решение с решением для струны при тех же
начальных условиях. Если заглянуть в лекцию 20, там можно найти решение
соответствующей задачи для струны:

\[
0 \le x \le \pi \qquad n = 0, 1, 2, \dots
\]

.

\pandocbounded{\includegraphics[keepaspectratio]{images/_page_158_Figure_1.jpeg}}

РИС 21.2. Простые колебания свободно опирающейся балки.

Другими словами, балка колеблется на более высоких частотах, чем струна.
Отметим, что в том и другом случаях высшие частоты кратны основным
частотам колебаний.

\section{ЗАМЕЧАНИЯ}\label{ux437ux430ux43cux435ux447ux430ux43dux438ux44f}

\begin{itemize}
\item
  \begin{enumerate}
  \def\labelenumi{\arabic{enumi}.}
  \tightlist
  \item
    Вообще говоря, балку можно закрепить тремя способами:
  \end{enumerate}

  \begin{itemize}
  \tightlist
  \item
    а) оставить свободной (незакрепленной),
  \item
    б) свободно опереть.
  \item
    в) жестко закрепить.
  \end{itemize}
\end{itemize}

Схематически эти способы закрепления изображены на рис. 21.3.

\begin{enumerate}
\def\labelenumi{\arabic{enumi}.}
\setcounter{enumi}{1}
\tightlist
\item
  Еще одна важная задача --- колебания балки с заделанным концом (см.
  рис. 213). Благодаря нестандартным граничным условиям
\end{enumerate}

\(u_x\)(0, t) = 0, \(u_{xx}\)(1, t) = 0, \(u_{xx}\)(1, t) = 0, решение
уже не будет выражаться в виде ряда из произведений синусов и косинусов.
Оно будет иметь более сложную структуру

\[
(\text{УЧП}) \quad u_{tt} + u_{xxxx} = 0, \quad 0 < x < 1, \quad 0 < t < \infty,
\]

где собственные функции \(X_n\)(x) являются линейными комбинациями
синусов, косинусов, гиперболических синусов и гиперболических косинусов.
Решение этой задачи можно найта в руководствах по теории упругости.

\pandocbounded{\includegraphics[keepaspectratio]{images/_page_159_Figure_2.jpeg}}

РИС. 21.3. Типичные задачи для балки: a---защемленная балка; b---балка,
жестко заделанная с обоих концов; a---левый конец жестко заделан, правый
свободно опирается.

\subsubsection{ЗАДАЧИ}\label{ux437ux430ux434ux430ux447ux438}

\begin{enumerate}
\def\labelenumi{\arabic{enumi}.}
\tightlist
\item
  Решите задачу для балки с закрепленным концом.
\end{enumerate}

\[
(\text{ГУ}) \quad
\begin{cases} u(0, t) &= 0, \\ u_x(0, t) &= 0, \\ u_{xx}(1, t) &= 0, \\ u_{xxx}(1, t) &= 0, \end{cases}
\]

\begin{verbatim}
\begin{cases}
\end{verbatim}

\end{document}
